\langsection[subsection]{Carrying Passengers and Cargo}{Перевозка пассажиров и грузов}


\begin{langblock}{0.425}
	\begin{eng}
	
	Only experienced riders should carry passengers or large loads. The extra weight changes the way the motorcycle handles, balances, turns, speeds up, and slows down. Before taking a passenger or heavy load on the street, practice away from traffic.
	\end{eng}
	\begin{rus}
		
		Неопытным, начинающим мотоциклистам не следует перевозить пассажиров или крупногабаритные грузы. Дополнительный вес влияет на управляемость мотоцикла, его устойчивость, повороты, разгон и торможение. Прежде чем выезжать на улицу с пассажиром или тяжелым грузом, потренируйтесь вдали от отживленного трафика.
\end{rus}
\end{langblock}



\langsection[subsubsection]{Equipment}{Оборудование}


\begin{langblock}{0.425}
	\begin{eng}
	
	To carry passengers safely:
	\begin{itemize}[noitemsep, nosep]
		\item Equip and adjust your motorcycle to carry passengers.
		\item Instruct the passenger before you start.
		\item Adjust your riding technique for the added weight.  
\end{itemize}	

	\medskip
	Equipment should include:
	\begin{itemize}[noitemsep, nosep]
		\item A proper seat -- large enough to hold both of you without crowding. You should not sit any farther forward than you usually do.
		\item Footrests and Handholds -- PA law requires that any motorcycle carrying a passenger, other than in a side car or enclosed cab, must be equipped with footrests and hand- holds for the passenger.
		\item Protective equipment -- the same protective gear recommended for operators.
	\end{itemize}
	
	Adjust the suspension to handle the additional weight. You will probably need to add a few pounds of pressure to the tires if you carry a passenger. (Check your owner's manual for appropriate settings.) While your passenger sits on the seat with you, adjust the mirrors and headlight according to the change in the motorcycle's angle.
	\end{eng}
	\begin{rus}
		
		Для безопасной перевозки пассажиров:
		\begin{itemize}[noitemsep, nosep]
			\item Оборудуйте и настройте свой мотоцикл
			\item Проинструктируйте пассажира перед началом поездки
			\item Адаптируйте манеру вождения с учетом дополнительного веса
	\end{itemize}	
	
	
	\medskip
	Оборудование должно включать:
	\begin{itemize}[noitemsep, nosep]
			\item Подходящее сиденье "--- достаточно большое, чтобы вы оба могли разместиться на нем, не теснясь. Вы не должны сдвигаться вперед больше, чем обычно.
			\item Подножки и поручни "--- законодательство штата Пенсильвания требует, чтобы любой мотоцикл, перевозящий пассажира (за исключением мотоциклов с коляской или закрытой кабиной), был оснащен подножками и поручнями для пассажира.
			\item Защитное снаряжение "--- для пассажира рекомендуется то же, что и для мотоциклиста.
		\end{itemize}
		
		Отрегулируйте подвеску с учетом дополнительного веса. Вероятно, вам придется увеличить давление в шинах. (Проверьте соответствующие таблицы в руководстве по эксплуатации.) Когда пассажир сидит с вами на сиденье, отрегулируйте зеркала и фару с учетом изменения угла наклона мотоцикла.
\end{rus}
\end{langblock}


\langsection[subsubsection]{Instructing Passengers}{Инструктаж}




\begin{langblock}{0.425}
	\begin{eng}
	
	Even if your passenger is a motorcycle rider, provide complete instructions before you start. Tell your passenger to:	
	\begin{itemize}[noitemsep, nosep]
		\item Get on the motorcycle only after you have started the engine.
		\item Sit as far forward as possible without crowding you.
		\item Hold firmly to your waist, hips, or belt.
		\item Keep both feet on the pegs, even when stopped.
		\item Keep legs away from the muffler(s), chains or moving parts.
		\item Stay directly behind you, leaning as you lean.
		\item Avoid unnecessary talk or motion.  
	\end{itemize}	
	
	\medskip
	Also, tell your passenger to tighten his or her hold when you:
	\begin{itemize}[noitemsep, nosep]
		\item Approach surface problems.
		\item Are about to start from a stop.
		\item Warn that you will make a sudden move.
	\end{itemize}
	\end{eng}
	\begin{rus}
		
		Даже если ваш пассажир "--- мотоциклист, все равно проговорите подробные инструкции. Попросите пассажира:
		
		\begin{itemize}[noitemsep, nosep]
			\item Садиться на мотоцикл только после запуска двигателя.
			\item Садиться как можно ближе вперед, но так, чтобы не было тесно.
			\item Крепко держаться за талию, бедра или ремень безопасности.
			\item Держать обе ноги на подножках, даже когда мотоцикл стоит.
			\item Держать ноги подальше от глушителя, цепи и движущихся частей.
			\item Держаться прямо за вами, наклоняясь в ту же сторону, что и вы.
			\item Избегать ненужных разговоров и телодвижений. 
		\end{itemize}	
		
		\medskip
		Кроме того, попросите пассажира крепче держаться, когда вы:
		\begin{itemize}[noitemsep, nosep]
			\item Подъезжаете к участку с неровностями.
			\item Собираетесь трогаться с места.
			\item Предупреждаете, что собираетесь совершить резкое движение.
		\end{itemize}
\end{rus}
\end{langblock}



\langsection[subsubsection]{Riding with Passengers}{Езда с пассажиром}


\begin{langblock}{0.425}
	\begin{eng}
	
	Your motorcycle will respond more slowly with a passenger on board. The heavier your passenger, the longer it may take to slow down, speed up, or turn --- especially on a light motorcycle.
	
	\begin{itemize}[noitemsep]
		\item Ride a little slower, especially when taking curves, corners, or bumps.
		\item Start slowing earlier as you approach a stop.
		\item Open up a larger cushion of space ahead and to the sides.
		\item Wait for larger gaps to cross, enter, or merge in traffic.
	\end{itemize}
	
	
	Warn your passenger of special conditions --- when you will pull out, stop quickly, turn sharply, or ride over a bump. Turn your head slightly to make yourself understood, but keep your eyes on the road ahead.
	\end{eng}
	\begin{rus}
		
		С пассажиром на борту ваш мотоцикл будет реагировать медленнее. Чем тяжелее пассажир, тем дольше может занять торможение, разгон или поворот "--- особенно на легком мотоцикле.
		
		\begin{itemize}[noitemsep]
			\item Езжайте медленнее, особенно при прохождении поворотов, изгибов или неровностей.
			\item Начинайте тормозить раньше при приближении к месту остановки.
			\item Оставляйте больше свободного пространства впереди и по сторонам.
			\item Дожидайтесь более широких интервалов, чтобы вклиниться в поток или пересечь дорогу.
		\end{itemize}
		
		
		Предупреждайте пассажира об особых ситуациях "--- когда вы собираетесь выехать, резко остановиться, совершить крутой поворот или проехать по неровности. Слегка поверните голову, чтобы пассажир вас точно понял, но не отрывайте взгляд от дороги.
\end{rus}
\end{langblock}

\langsection[subsubsection]{Carrying Loads}{Перевозка груза}



\begin{langblock}{0.425}
	\begin{eng}
	
	Most motorcycles are not designed to carry much cargo. Small loads can be carried safely if positioned and fastened properly.
	
	
	\begin{itemize}[noitemsep]
		\item Keep the Load Low -- Fasten loads securely, or put them in saddle bags.  Piling loads against a sissybar or frame on the back of the seat raises the motorcycle's center of gravity and disturbs its balance.
		\item Keep the Load Forward -- Place the load over, or in front of, the rear axle. Tank bags keep loads forward, but use caution when loading hard or sharp objects. Make sure tank bag does not interfere with handlebars or controls. Mounting loads behind the rear axle can affect how the motorcycle turns and brakes. It can also cause a wobble.
		\item Distribute the Load Evenly -- Load saddlebags with about the same weight. An uneven load can cause the motorcycle to drift to one side.
		\item Secure the Load -- Fasten the load securely with elastic cords (bungee cords or nets). Elastic cords with more than one attachment point per side are more secure. A tight load won't catch in the wheel or chain, causing it to lock up and skid. Rope tends to stretch and knots come loose, permitting the load to shift or fall.
		\item Check the Load -- Stop and check the load every so often to make sure it has not worked loose or moved.
	\end{itemize}
	\end{eng}
	\begin{rus}
		
		Большинство мотоциклов не рассчитаны на перевозку грузов. Небольшой багаж можно безопасно перевозить при условии его правильного размещения и крепления.
		
		\begin{itemize}[noitemsep]
			\item Размещайте груз ниже "--- надежно закрепляйте или помещайте их в сумки/кофры. Укладка грузов у спинки сиденья или рамы сзади повышает центр тяжести мотоцикла и нарушает его баланс.
			\item Располагайте груз впереди "--- размещайте груз над задней осью или перед ней. Сумки на баке позволяют размещать груз впереди, но будьте осторожны с твердыми и острыми предметами. Убедитесь, что сумка на баке не мешает рулю или органам управления. Размещение груза за задней осью может повлиять на поведение мотоцикла при поворотах и торможении. Это также может вызвать биение переднего колеса (Шимми/Wobble).
			\item Равномерно распределите груз "--- загружайте боковые сумки равномерно, чтобы избежать сноса мотоцикла в сторону.
			\item Закрепите груз "--- надежно закрепите груз с помощью эластичных шнуров (банджи-шнуров или сеток). Эластичные шнуры с несколькими точками крепления обеспечивают большую надежность. Плотно закрепленный груз не зацепится за колесо или цепь, что могло бы привести к блокировке и заносу. Веревка имеет свойство растягиваться, а узлы развязываться.
			\item Проверяйте груз "--- время от времени останавливайтесь и проверяйте груз, чтобы убедиться, что он крепеж не ослаб.
		\end{itemize}
\end{rus}
\end{langblock}









